\scanpage{30}
\mspage{28}
% archive sheet number: 29
% hand: Grothendieck
% sheet_kind: text
% element: text | heading | 6. Cas d'un corps résiduel de type fini.
6. \uline{Cas d'un corps résiduel de type fini.}

% element: text | upper half, lines 2-10 | Proposition 6.1. Soit pi un groupe profini
\uline{Proposition} 6.1. Soit $\pi$ un groupe profini, extension d'un groupe profini $\pi_0$ par un groupe profini \deleted{\unsure{commutatif}} \inserted{résoluble} $I$, \deleted{\unsure{de sorte que|tel que} $\pi_0$ \unsure{opère sur} $I$, \unsure{le}} supposons que pour \unsure{tt|$\forall$} ss-groupe ouvert $\pi'_0$ de $\pi_0$, \deleted{le groupe} \inserted{[N.B. si $I$ commutatif, \unsure{celà|cela} signifie $I'_{\pi'_0}=\varprojlim (I'_\alpha)_{\pi'_0}$ \deleted{$\ne [I,\pi']$} \unsure{[}(où $I_\alpha$ parcourt les quotients \deleted{finis} de $I$ par des sous-groupes ouverts, invariants dans $\pi$\unsure{]|)}} \deleted{est un groupe \unsure{où} $\pi'$ \unsure{est} l'image inverse de $\pi'_0$ dans $\pi$} fini] \deleted{que} posant $I'=\pi'\cap I$, $[I',\pi']$ est d'indice fini dans $\unsure{I|I'}$, \unsure{\{où|(où} $\pi'$ est l'image inverse de $\pi'_0$ dans $\pi$\unsure{\}|)}. Alors pour tte représentation $\ell$-adique de $\pi$, la représentation induite sur $I$ est quasi-unipotente.

% element: margin note | left margin beside statement, slanted | Cf. nouvelle version
% slanted
\margin{Cf. nouvelle version à la fin du §}

% element: text | middle, lines 13-18 | Soit E un vectoriel de dim finie
Soit $E$ un vectoriel de dim finie sur $\mathbb Q_\ell$ \deleted{dans} lequel on fait opérer $\pi$. Soit \deleted{$\unsure{G\text{ et }(|G}$} $G$ ($H$) la composante connexe du groupe algébrique défini par \inserted{enveloppe de} $\pi$ (resp. $I$), de sorte que $H\subset G$, $H$ \inserted{\deleted{commutatif} résoluble,} invariant dans $G$. L'hypothèse \deleted{signifie} \inserted{implique} que $[G,H]=H$.

% element: text | lower half, lines 19-25 | On a encore H=T.U
On a \unsure{encore} $H=\deleted{H_r\unsure{\ltimes}H_u}\inserted{T.U}$ (\inserted{produit semi-direct de} partie \deleted{\unsure{réductive}}\inserted{torique} et partie unipotente), et $G$ opère sur \deleted{$H_r$ et $H_u$ séparément}, \inserted{$U$, donc sur $H/U\simeq T$} et opère trivialement sur $\deleted{H_r}\inserted{T}$ (\deleted{qui est un tore} $G$ étant connexe). Donc $[G,H]=H$ implique \deleted{$H_r=1$}, \inserted{$[G,T]=T$ donc $T=1$,} donc $H$ est unipotent, cqfd.
